\section{Beyond Transformers：JEPA、world models 与 state-space models}

课程最后没有宣布 Transformer 时代结束，而是提出两类替代或补充方向：world model/JEPA 改变预测对象，从 token 转向抽象 latent state；state-space model（SSM）改变序列计算，从显式全局 attention 转向递归状态更新。二者分别针对表示目标与计算复杂度。

\subsection{为什么要越过 next-token prediction}

Transformer 在多领域占据主导，但仍面临效率、长上下文执行和结构化 world understanding。slide 151 将 emerging directions 分为 world models 与 SSMs。这里的“beyond”应读作扩展设计空间，而不是否定 attention：很多实际系统会把 Transformer、SSM、retrieval 与 simulator 混合。

\lecturefigure{slide-151.jpg}{Beyond Transformers：效率、长上下文与 world understanding 的缺口}{CS25 V6 official deck, p.~151}

\teachervoice{01:12:51--01:13:20，老师说 Transformer 仍主导，但在效率、long-context reasoning 与 world understanding 上存在限制。课堂提示：下一代系统很可能是组合，而非单一架构完全取代。}

\subsection{World models 与 JEPA：预测 latent，而非重建全部细节}

Joint Embedding Predictive Architecture（JEPA）在 latent space 预测目标表示。设 context encoder 为 $f_{\theta}$、target encoder 为 $f_{\xi}$、predictor 为 $g_{\phi}$，则可优化
\[
\mathcal{L}_{\mathrm{JEPA}}
=\left\|g_{\phi}(f_{\theta}(x_{\mathrm{context}}))-\operatorname{sg}\big(f_{\xi}(x_{\mathrm{target}})\big)\right\|_2^2.
\]
$\operatorname{sg}$ 表示 stop-gradient。目标是预测与未来/缺失区域相关的抽象表示，而不是逐像素重建所有不可预测细节。难点是防止 representation collapse，并确保 latent 保留 planning 所需状态。

\lecturefigure{slide-152.jpg}{World models 与 JEPA：从 next-token prediction 转向结构化环境表示}{CS25 V6 official deck, p.~152}

\lecturefigure{slide-153.jpg}{JEPA 机制：context encoder、latent prediction 与抽象目标}{CS25 V6 official deck, p.~153}

\teachervoice{01:13:02--01:14:38，老师把 JEPA 描述为学习环境结构和未来状态的方向，而不是“预测更多像素”。讲义提醒：latent prediction 是否形成可用于规划的 world model 仍需任务证据。}

\begin{warningbox}{World model 不能只靠名字成立}
一个模型能预测视频 latent、下一个 observation 或 token，不等于它拥有因果、可干预、可规划的世界模型。至少要检验 state sufficiency、counterfactual consistency、long-horizon rollout 和 action-conditioned prediction。
\end{warningbox}

\subsection{SSM：用状态更新换掉显式平方交互}

State Space Model 将序列写成状态更新：
\[
h_t=\bar{A}_t h_{t-1}+\bar{B}_t x_t,
\qquad y_t=C_t h_t+D x_t.
\]
$h_t$ 是压缩状态，$\bar{A}_t,\bar{B}_t,C_t$ 可固定或由输入选择。Selective SSM（如 Mamba）让部分参数依赖输入，以决定写入、保留和读取。训练可借助 scan/convolution 技巧并行，推理状态通常按序列长度线性更新。

\lecturefigure{slide-154.jpg}{State Space Model：连续状态、线性时间与 Mamba 示例}{CS25 V6 official deck, p.~154}

\lecturefigure{slide-155.jpg}{SSM tradeoffs：效率、长序列与 attention 灵活性的对照}{CS25 V6 official deck, p.~155}

\teachervoice{01:14:38--01:16:20，老师强调 SSM 在长序列和线性复杂度上有吸引力，但 attention 在某些需要灵活内容寻址的任务上更强，当前结论仍未封闭。}

\begin{knowledgebox}{术语消化：attention 与 SSM 的核心差异}
Attention 为每个 query 显式访问整段 context，内容寻址灵活但 score matrix 昂贵；SSM 把过去压进固定状态，推理高效但可能丢失可随机访问细节。Hybrid architecture 会在局部/关键层使用 attention，在大部分层使用 state-space 或 convolution，以寻找 workload-specific Pareto point。
\end{knowledgebox}

比较 JEPA、SSM 与 Transformer 时还要固定 representation 和任务。若 JEPA 使用更强视觉 encoder，而 baseline 直接预测像素，性能差异不全来自 objective；若 SSM 使用更长 context 或不同 tokenizer，吞吐差异也不能直接归因于状态更新。公平实验至少匹配参数、训练 token、FLOPs、wall-clock、序列长度和数据增强，并分别报告 prefill、decode、memory 与质量。更重要的是测试失败形态：attention 是否在精确复制和任意位置检索上更强，SSM 是否在流式长序列更稳，JEPA latent 是否真正支持 action-conditioned planning。只有这些 workload-level 结果才能决定组合架构，而不是由复杂度符号替代实测。

\subsection{本章小结}

JEPA/world models 改变学习目标，SSM 改变计算路径。前者追求对环境状态更抽象、可预测的表示；后者追求长序列效率。两者都不是凭概念名自动优于 Transformer，必须在明确 tokenization、数据、任务、compute 与评估下比较。

